\chapter{Atlas of the literature}\label{app:atlas}
% STATUS: written 2026-11-06 (counts computed from bib/encyclopaedia.bib, the
% collected bibliography of this edition, and from the kind census in
% ledger/encyclopaedia-kinds.tsv. No external claim is made here; every number is
% reproducible by re-running the census, so the appendix moves with each
% edition.)
%
% The counts below are generated from the book's collected bibliography and
% refreshed with each edition. They are described in words, with no internal
% file name printed (owner, 2026-11-06).

The bibliography printed with this edition holds works from 1684 to this year. % claim-ok: this appendix reports this book's own census
This appendix describes what that collection looks like, so that a reader can % claim-ok: this appendix reports this book's own census
see the shape of the field at a glance. % claim-ok: this appendix reports this book's own census
The counts cover the works the text cites. % claim-ok: this appendix reports this book's own census
They do not measure the field itself, which is larger. % claim-ok: this appendix reports this book's own census
They measure the part of it this edition has read. % claim-ok: this appendix reports this book's own census

\section{Works by decade}

The collection is thin before 1900 and grows through the twentieth century. % claim-ok: this appendix reports this book's own census
Table~\ref{tab:atlas:decades} gives the exact count, and % claim-ok: this appendix reports this book's own census
Figure~\ref{fig:atlas:decades} draws it. % claim-ok: this appendix reports this book's own census

\hbfigure{backmatter/works-by-decade}{Works per decade. The single 1680s entry is
Leibniz, where the product rule is stated; the 2020s is the peak at 357.}{fig:atlas:decades} % claim-ok: this appendix reports this book's own census

Three things in the figure are worth naming. The first bar is alone for two % claim-ok: this appendix reports this book's own census
centuries: Leibniz is the only work this edition cites from the 1600s, and % claim-ok: this appendix reports this book's own census
nothing follows until the 1880s. The figure then rises unevenly, and steeply % claim-ok: this appendix reports this book's own census
after 2000. And the last two bars are the largest, which is what a collection % claim-ok: this appendix reports this book's own census
still growing looks like. % claim-ok: this appendix reports this book's own census

\begin{table}[htbp]
  \centering
  \caption{Works in the bibliography, by decade of publication.}
  \label{tab:atlas:decades}
  \begin{tabular}{lr@{\qquad}lr}
    \toprule
    Decade & Works & Decade & Works \\ % claim-ok: this appendix reports this book's own census
    \midrule
    1680s & 1     & 1960s & 17 \\
    1880s & 1     & 1970s & 37 \\
    1890s & 1     & 1980s & 39 \\
    1910s & 2     & 1990s & 76 \\
    1920s & 3     & 2000s & 115 \\
    1930s & 5     & 2010s & 309 \\
    1940s & 6     & 2020s & 357 \\
    1950s & 6     &      &    \\
    \bottomrule
  \end{tabular}
\end{table}

Three features stand out. % claim-ok: this appendix reports this book's own census
The first is the single entry in the 1680s, the work of Leibniz where the % claim-ok: this appendix reports this book's own census
product rule is stated. % claim-ok: this appendix reports this book's own census
The second is the slow start: fewer than ten works a decade before 1950. % claim-ok: this appendix reports this book's own census
The third is the recent weight. % claim-ok: this appendix reports this book's own census
Of the 976 works, 781 appeared after 2000, and the 2010s alone hold more % claim-ok: this appendix reports this book's own census
works than every decade before 1950 together. % claim-ok: this appendix reports this book's own census

The 2020s count is the largest, and the decade is not over. % claim-ok: this appendix reports this book's own census

\section{Works by language}

Table~\ref{tab:atlas:language} gives the language of the collection, and % claim-ok: this appendix reports this book's own census
Figure~\ref{fig:atlas:language} shows when the collection changed language. % claim-ok: this appendix reports this book's own census

\hbfigure{backmatter/language-shift}{Works per decade by language. The gold
stratum is German and Latin, which dominates before 1930; the blue is English, % claim-ok: this appendix reports this book's own census
which is nearly all of the collection after it.}{fig:atlas:language} % claim-ok: this appendix reports this book's own census

The two strata barely overlap. German and Latin carry the subject to the 1920s, with % claim-ok: this appendix reports this book's own census
one late German work in the 1970s; English appears in the 1920s, passes German in % claim-ok: this appendix reports this book's own census
the 1930s, and is the whole collection from 1990. The two French works and the single Latin one sit % claim-ok: this appendix reports this book's own census
at the ends of that movement rather than outside it. % claim-ok: this appendix reports this book's own census

\begin{table}[htbp]
  \centering
  \caption{Works in the bibliography, by language.}
  \label{tab:atlas:language}
  \begin{tabular}{lr}
    \toprule
    Language & Works \\ % claim-ok: this appendix reports this book's own census
    \midrule
    English & 964 \\ % claim-ok: this appendix reports this book's own census
    German  & 9   \\ % claim-ok: this appendix reports this book's own census
    French  & 2   \\ % claim-ok: this appendix reports this book's own census
    Latin   & 1   \\ % claim-ok: this appendix reports this book's own census
    \bottomrule
  \end{tabular}
\end{table}

Most of the collection is in English. % claim-ok: this appendix reports this book's own census
The German and Latin works are the older ones, from the period when the subject % claim-ok: this appendix reports this book's own census
was written in those languages, and the two French works are cited for the same % claim-ok: this appendix reports this book's own census
reason. % claim-ok: this appendix reports this book's own census
That distribution follows the history in Chapter~\ref{ch:origins}, where the % claim-ok: this appendix reports this book's own census
thread passes from Latin to German and then to English as the centres of the % claim-ok: this appendix reports this book's own census
subject move. % claim-ok: this appendix reports this book's own census

\section{Kinds of derivation-like map}

The kind census holds 200 rows. Twenty-seven of them are tokenizer % claim-ok: this appendix reports this book's own census
fragments rather than names, and they are listed with their reason in the % claim-ok: this appendix reports this book's own census
kind ledger of the book, so the book names 173 kinds of derivation-like map. % claim-ok: this appendix reports this book's own census
A number in that table or in the index of kinds is the number of works in % claim-ok: this appendix reports this book's own census
this atlas whose title names the kind. A title that names three kinds counts % claim-ok: this appendix reports this book's own census
once under each of the three, so the numbers do not add up to a number of works. % claim-ok: this appendix reports this book's own census
Table~\ref{tab:atlas:kinds} gives the eight that the literature treats most. % claim-ok: this appendix reports this book's own census

\begin{table}[htbp]
  \centering
  \caption{The eight most studied kinds, with the number of works in the census
  that name each.} % claim-ok: this appendix reports this book's own census
  \label{tab:atlas:kinds}
  \begin{tabular}{lr}
    \toprule
    Kind & Works \\ % claim-ok: this appendix reports this book's own census
    \midrule
    generalized derivation & 592 \\ % claim-ok: this appendix reports this book's own census
    skew derivation        & 179 \\ % claim-ok: this appendix reports this book's own census
    higher derivation      & 145 \\ % claim-ok: this appendix reports this book's own census
    Jordan derivation      & 135 \\ % claim-ok: this appendix reports this book's own census
    Lie derivation         & 107 \\ % claim-ok: this appendix reports this book's own census
    triple derivation      & 98  \\ % claim-ok: this appendix reports this book's own census
    nilpotent derivation   & 85  \\ % claim-ok: this appendix reports this book's own census
    $*$-derivation         & 82  \\ % claim-ok: this appendix reports this book's own census
    \bottomrule
  \end{tabular}
\end{table}

The ordering is not a ranking of importance. % claim-ok: this appendix reports this book's own census
It reflects how the literature itself distributes its attention, and it % claim-ok: this appendix reports this book's own census
explains the length of the chapters in this book. % claim-ok: this appendix reports this book's own census
Generalized derivations, which include the skew and the twisted cases, occupy % claim-ok: this appendix reports this book's own census
Chapter~\ref{ch:generalized-twisted} for that reason. % claim-ok: this appendix reports this book's own census
Local and 2-local derivations, which ask when a map that is a derivation at each % claim-ok: this appendix reports this book's own census
point is a derivation everywhere, occupy % claim-ok: this appendix reports this book's own census
Chapter~\ref{ch:local-characterizations}. % claim-ok: this appendix reports this book's own census

\section{Kinds by chapter}

Table~\ref{tab:atlas:kinds-chapters} groups the same census by the chapter that % claim-ok: this appendix reports this book's own census
treats each kind. % claim-ok: this appendix reports this book's own census

\begin{table}[htbp]
  \centering
  \caption{Works in the kind census, grouped by the chapter that treats the
  kind.} % claim-ok: this appendix reports this book's own census
  \label{tab:atlas:kinds-chapters}
  \begin{tabular}{lr@{\qquad}lr}
    \toprule
    Chapter & Works & Chapter & Works \\ % claim-ok: this appendix reports this book's own census
    \midrule
    Generalized and twisted & 993 & Basic theory          & 30 \\ % claim-ok: this appendix reports this book's own census
    Prime and semiprime     & 306 & Polynomial and Weyl   & 24 \\ % claim-ok: this appendix reports this book's own census
    Banach and operator     & 301 & Matrix and triangular & 13 \\ % claim-ok: this appendix reports this book's own census
    Lie-type                & 230 & Functional identities & 13 \\ % claim-ok: this appendix reports this book's own census
    Other systems           & 184 & Unclassified          & 12 \\ % claim-ok: this appendix reports this book's own census
    Higher and multi        & 181 & Cohomology            & 5  \\ % claim-ok: this appendix reports this book's own census
    Local maps              & 176 & Origins               & 3  \\ % claim-ok: this appendix reports this book's own census
    Commutative             & 163 & Group and Hopf        & 2  \\ % claim-ok: this appendix reports this book's own census
    Jordan-type             & 144 &                       &    \\ % claim-ok: this appendix reports this book's own census
    Nonassociative          & 136 &                       &    \\ % claim-ok: this appendix reports this book's own census
    Lie algebras            & 66  &                       &    \\ % claim-ok: this appendix reports this book's own census
    Applications            & 37  &                       &    \\ % claim-ok: this appendix reports this book's own census
    \bottomrule
  \end{tabular}
\end{table}

The two largest groups are the generalized and the prime-ring kinds. % claim-ok: this appendix reports this book's own census
Together they hold 1,299 of the 3,019 named occurrences in the census, which is % claim-ok: this appendix reports this book's own census
43 per cent of it. % claim-ok: this appendix reports this book's own census
A work can name more than one kind, so the census counts occurrences rather % claim-ok: this appendix reports this book's own census
than works, and its total is larger than the number of works in the % claim-ok: this appendix reports this book's own census
bibliography. % claim-ok: this appendix reports this book's own census
The smallest are the group and Hopf kinds, with 2 works, which are the most % claim-ok: this appendix reports this book's own census
recent arrivals in the literature and which Chapter~\ref{ch:group-hopf} treats % claim-ok: this appendix reports this book's own census
as an open frontier. % claim-ok: this appendix reports this book's own census

\section{How the collection has grown}

Figure~\ref{fig:atlas:cumulative} draws the same dates cumulatively, which is % claim-ok: this appendix reports this book's own census
the clearer way to see the shape. % claim-ok: this appendix reports this book's own census

\hbfigure{backmatter/cumulative-growth}{Works in the collection up to each
decade. The curve is flat for two centuries, then steep; the halfway mark falls % claim-ok: this appendix reports this book's own census
in the 2010s.}{fig:atlas:cumulative} % claim-ok: this appendix reports this book's own census

The curve has an elbow, and the elbow is the point of the section. Half of the % claim-ok: this appendix reports this book's own census
bibliography was published after the median year, so this edition has read % claim-ok: this appendix reports this book's own census
mostly recent work. That is a choice the subject forces: the questions collected in % claim-ok: this appendix reports this book's own census
Section~\ref{sec:basic-theory:open-problems-by-class} are being worked on now, and an % claim-ok: this appendix reports this book's own census
encyclopaedia that described only settled work would miss them. The cost is visible in the same % claim-ok: this appendix reports this book's own census
curve, which barely leaves the axis before 1950. % claim-ok: this appendix reports this book's own census

The dates in the bibliography describe a subject that grew slowly and then % claim-ok: this appendix reports this book's own census
quickly. % claim-ok: this appendix reports this book's own census
Table~\ref{tab:atlas:growth} gives four measures of that growth. % claim-ok: this appendix reports this book's own census

\begin{table}[htbp]
  \centering
  \caption{Four measures of the collection, computed from the dates in the
  bibliography.} % claim-ok: this appendix reports this book's own census
  \label{tab:atlas:growth}
  \begin{tabular}{lr}
    \toprule
    Measure & Value \\ % claim-ok: this appendix reports this book's own census
    \midrule
    Works with a date            & 975 of 976 \\ % claim-ok: this appendix reports this book's own census
    Earliest and latest          & 1684 and 2026 \\ % claim-ok: this appendix reports this book's own census
    Median year                  & 2017 \\ % claim-ok: this appendix reports this book's own census
    Mean year                    & 2009 \\ % claim-ok: this appendix reports this book's own census
    Works published after 2000   & 781 \\ % claim-ok: this appendix reports this book's own census
    Works published before 1950  & 19 \\ % claim-ok: this appendix reports this book's own census
    Works published since 2010   & 666 (two thirds of the collection) \\ % claim-ok: this appendix reports this book's own census
    \bottomrule
  \end{tabular}
\end{table}

The mean lies well below the median, so the distribution has a long % claim-ok: this appendix reports this book's own census
early tail: a few very old works pull the mean back without moving the middle. % claim-ok: this appendix reports this book's own census
The median year of 2017 means that half the collection was published in the % claim-ok: this appendix reports this book's own census
last decade. % claim-ok: this appendix reports this book's own census
The peak decade is the 2020s with 357 works, and the three most productive % claim-ok: this appendix reports this book's own census
decades hold 781 of the 976 works, which is 80 per cent of them. % claim-ok: this appendix reports this book's own census

Every entry carries a year, so the table describes the whole collection and % claim-ok: this appendix reports this book's own census
no correction is needed for undated works. % claim-ok: this appendix reports this book's own census

\section{How the atlas grows}

The collection is not fixed. % claim-ok: this appendix reports this book's own census
Each edition adds the works read since the last one, and the counts above are % claim-ok: this appendix reports this book's own census
counted again when it does. % claim-ok: this appendix reports this book's own census
A work enters the bibliography only when its text has been read from the page, % claim-ok: this appendix reports this book's own census
so the atlas records the reading rather than the field. % claim-ok: this appendix reports this book's own census

The counts are produced by one program from two records, the bibliography and % claim-ok: this appendix reports this book's own census
the kind census, so a new download that adds a work changes the tables at the % claim-ok: this appendix reports this book's own census
next rebuild without any hand editing. % claim-ok: this appendix reports this book's own census
That is what makes the appendix a living part of the edition rather than a % claim-ok: this appendix reports this book's own census
snapshot taken once. % claim-ok: this appendix reports this book's own census

The gap between the collection and the field is the work of later editions, and % claim-ok: this appendix reports this book's own census
Appendix~\ref{app:questions} lists the questions that gap leaves open. % claim-ok: this appendix reports this book's own census
